\hnsection{THE LIE GROUPS-LIE ALGEBRAS ``DICTIONARY''}

The theory of ``finite continuous'' groups, developed by Lie in numerous memoirs beginning in 1874, is expounded systematically in the imposing treatise ``Theorie der Transformationsgruppen'' ({[}4{]}, 1888--1893{]}), written in collaboration with F. Engel§; it is the object of study in the first volume and the last five chapters of the third, the second being devoted to contact transformations.

As the title indicates, this work is only concerned with transformation groups in the sense of equations (4), where the space of ``variables'' \(x_{i}\) and the space of ``parameters'' \(a_{i}\) play initially such important roles. On the other hand, the concept of ``abstract'' group had not been clearly isolated at that period; when in 1885 ({[}3/j{]}, §5) Lie noted that in the notation of (10) the equation \(w = \mathrm{H}(u, v)\) which gives the parameters of the composition of two transformations of the group defines a new group, he considered it as a transformation group on the space of parameters, thus obtaining what he called the ``parameter group'' (he even obtained two, which are just the group of left translations and the group of right translations†).

The variables \(x_{i}\) and the parameters \(a_{i}\) in equations (4) were in principle assumed to be complex (except in Chapters XIX--XXIV of volume 3) and the functions \(f_{i}\) to be analytic; Lie and Engel were of course aware of the fact that these functions are not in general defined for all complex values of the \(x_{i}\) and \(a_{i}\) and that therefore the composition of such functions raises serious difficulties ({[}4{]}, vol.~1, pp.~15--17, pp.~33--40 and passim); and although throughout they almost always wrote as if the composition of the transformations they were studying was possible without restriction, this was no doubt for the convenience of the results and they explicitly reaffirmed the ``local'' point of view whenever necessary (cf.~loc. cit., p.~168 or 189 for example or ibid., vol.~3, p.~2, note at the bottom of the page); in other words, the mathematical object they studied is close to what we call in this treatise a law chunk of operation. They did not refrain, on occasions, from considering global groups, for example the 4 series of classical groups ({[}4{]}, vol.~3, p.~682), but do not appear to have asked themselves the question of what in general constitutes a ``global group''; they were content to obtain, for the ``parameters'' of the classical groups (the ``variables'' of these groups introduced no difficulty, since the transformations in question are linear transformations of \(\mathbf{C}^n\) ) ``local'' systems of parameters in a neighbourhood of the identity transformation, without worrying about the domain of validity of the formulae they were writing down. They however set themselves a problem which arises neatly out of the local theory: the study of ``mixed'' groups, that is groups with a finite number of connected components, such as the orthogonal group ({[}4{]}, vol.~1, p.~7). They presented this study as a study of a set of transformations stable under composition and passage to the inverse which is the union of sets \(\mathrm{H}_j\) each of which is described by systems of functions ( \(f_i^{(D)}\) ) as in (4); the number of (essential) parameters of each \(\mathrm{H}_j\) is even a priori assumed to depend on \(j\) , but they showed that in fact this number is the same for all the \(\mathrm{H}_j\) . Their principal result was then the existence of a finite continuous group G such that \(\mathrm{H}_j = \mathrm{G} \circ h_j\) for some \(h_j \in \mathrm{H}_j\) and for all \(j\) ; they also established that G is normal in the mixed group and noted that the determination of the invariants of the latter reduces to that of the invariants of G and a discontinuous group ({[}4{]}, vol.~1, Chapter 18).

The general theory developed in {[}4{]} ended (without the authors saying so very systematically) by achieving a ``dictionary'' translating the properties of ``finite continuous'' groups into those of the set of their infinitesimal transformations. It is based on the ``three theorems of Lie'', each of which consists of an assertion and its converse.

The first theorem ({[}4{]}, vol.~1, pp.~33 and 72 and vol.~3, p.~563) affirms in the first place that if in (4) the parameters are effective, the functions \(f_{i}\) satisfy a system of partial differential equations of the form

\[
\frac {\partial f _ {i}}{\partial a _ {j}} = \sum_ {k = 1} ^ {r} \xi_ {k i} (f (x, a)) \psi_ {k j} (a) \quad (1 \leqslant i \leqslant n)\tag{15}
\]

where the matrix \((\xi_{kl})\) is of maximum rank and \(\det(\psi_{kj}) \neq 0\) ; conversely, if the functions \(f_{i}\) have this property, formulae (4) define a group germ of transformations.

The second theorem ({[}4{]}, vol.~1, pp.~149 and 158, and vol.~3, p.~590) gives relations between the \(\xi_{ki}\) on the one hand and the \(\psi_{ij}\) on the other: the conditions on the \(\xi_{ki}\) can be written in the form

\[
\sum_ {k = 1} ^ {n} \left(\xi_ {i k} \frac {\partial \xi_ {f l}}{\partial x _ {k}} - \xi_ {j k} \frac {\partial \xi_ {t l}}{\partial x _ {k}}\right) = \sum_ {k = 1} ^ {r} c _ {i j} ^ {k} \xi_ {k l} (1 \leqslant i, j \leqslant r, 1 \leqslant l \leqslant n)\tag{16}
\]

where the \(c_{ii}^{k}\) are constants \((1\leqslant i,j,k\leqslant r)\) skew-symmetric in \(i,j\) . The conditions on the \(\psi_{ij}\) , in the form given by Maurer {[}10{]}, are:

\[
\frac {\partial \psi_ {k l}}{\partial a _ {m}} - \frac {\partial \psi_ {k m}}{\partial a _ {l}} = \frac {1}{2} \sum_ {1 \leqslant i, j \leqslant r} c _ {i j} ^ {k} (\psi_ {i l} \psi_ {j m} - \psi_ {j l} \psi_ {i m}) (1 \leqslant k, l, m \leqslant r).\tag{17}
\]

By introducing the contragradient matrix \((\alpha_{ij})\) of \((\psi_{ij})\) and the infinitesimal transformations

\[
\mathbf {X} _ {k} = \sum_ {i = 1} ^ {n} \xi_ {k i} \frac {\partial}{\partial x _ {i}}, \quad \mathbf {A} _ {k} = \sum_ {j = 1} ^ {r} \alpha_ {k j} \frac {\partial}{\partial a _ {j}} (1 \leqslant k \leqslant r)\tag{18}
\]

\begin{enumerate}
\def\labelenumi{(\arabic{enumi})}
\setcounter{enumi}{15}
\item
  and (17) can be written respectively:
\item
\end{enumerate}

\[
[ \mathrm{X} _ {i}, \mathrm{X} _ {j} ] = \sum_ {k = 1} ^ {r} c _ {i j} ^ {k} \mathrm{X} _ {k} (1 \leqslant i, j \leqslant r).\tag{20}
\]

\[
[ \mathbf {A} _ {i}, \mathbf {A} _ {j} ] = \sum_ {k = 1} ^ {r} c _ {i j} ^ {k} \mathbf {A} _ {k}
\]

Conversely, if r infinitesimal transformations \(X_{k}\) ( \(1 \leqslant k \leqslant r\) ) are given which are linearly independent and satisfy conditions (19), the one-parameter subgroups generated by these transformations generate a transformation group in r essential parameters.

Finally, the third theorem ({[}4{]}, vol.~1, pp.~170 and 297 and vol.~3, p.~597) reduces the determination of the systems of infinitesimal transformations \((\mathbf{X}_k)_{1 \leqslant k \leqslant r}\) satisfying (19) to a purely algebraic problem: the following must hold:

\begin{enumerate}
\def\labelenumi{(\arabic{enumi})}
\setcounter{enumi}{20}
\tightlist
\item
\end{enumerate}

\[
c _ {i j} ^ {k} + c _ {j i} ^ {k} = 0\tag{22}
\]

\[
\sum_ {l = 1} ^ {r} \left(c _ {i l} ^ {m} c _ {j k} ^ {l} + c _ {k i} ^ {m} c _ {i j} ^ {l} + c _ {j l} ^ {m} c _ {k l} ^ {l}\right) = 0 \quad (1 \leqslant i, j, k, m \leqslant r).
\]

Conversely,† if (21) and (22) are satisfied, there exists a system of infinitesimal transformations satisfying relations (19), whence a transformation group with r parameters (in other words, the linear combinations of the \(X_{k}\) with constant coefficients form a Lie algebra and conversely every finite-dimensional Lie algebra can be obtained in this way).

These results are completed by studying the questions of isomorphisms. Two transformation groups are called similar if it is possible to pass from one to the other by an invertible coordinate transformation on the variables and an invertible coordinate transformation on the parameters: from the very beginning of his research, Lie had been concerned naturally with this notion when defining the ``canonical parameters''. He showed that two groups are similar if, by a transformation on the ``variables'', it is possible to transform the infinitesimal transformations of the one into those of the other ({[}4{]}, vol.~1, p.~329). A necessary condition for this to be the case, is that the Lie algebras of these two groups be isomorphic, which Lie expressed by saying that the groups are ``gleichzusammengesetzt''; but this condition is not sufficient and a whole chapter ({[}4{]}, vol.~1, Chapter 19) is devoted to obtaining supplementary conditions assuring that the groups are ``similar''. The theory of permutation groups on the other hand provided the notion of ``holohedric isomorphism'' of two such groups (isomorphism of the underlying ``abstract'' groups); Lie transposed this notion to transformation groups and showed that two such groups are ``holohedrically isomorphic'' if and only if their Lie algebras are isomorphic ({[}4{]}, vol.~1, p.~418). In particular, every transformation group is holohedrically isomorphic to each of its parameter groups and this shows that, when we wish to study the structure of the group, the ``variables'' on which it operates are of little importance and that in fact all that matters is the Lie algebra.†

Always by analogy with the theory of permutation groups, Lie introduced the notions of subgroups, normal subgroups and ``merihedric isomorphisms'' (surjective homomorphisms) and showed that they correspond to those of subalgebras, ideals and surjective Lie algebra homomorphisms; he had already previously come across an especially important example of a ``merihedric isomorphism'', the adjoint representation, and had recognized its relations with the centre of the group ({[}3f{]}, § 3, no. 9). For these results, as for the fundamental theorems, the essential tool is the Jacobi-Clebsch Theorem giving the complete integrability of a differential system (one of the forms of the theorem called ``{[}Frobenius's''); he however gave a new proof using one-parameter groups ({[}4{]}, vol.~1, Chapter 6).

The notions of transitivity and primitivity, so important for permutation groups, occurred just as naturally for ``finite continuous'' transformation groups and the Lie--Engel treatise made a detailed study of this ({[}4{]}, vol.~1, Chapter 13 and passim); the relations with the stabilizer subgroups of a point and the notion of homogeneous space were perceived (insofar as was possible without taking a global point of view) ({[}4{]}, vol.~1, p.~425).

Finally the ``dictionary'' was completed, in {[}4{]}, by the introduction of the notions of derived group and solvable group (called ``integrable group'' by Lie; this terminology, suggested by the theory of differential equations, remained in use until the works of H. Weyl) ({[}4{]}, vol.~1, p.~261 and vol.~3, pp.~678--679); the relation between commutators and brackets had elsewhere been perceived by Lie in 1883 ({[}3{]}, vol.~5, p.~358).

